\documentclass[12pt]{article}
\usepackage[ left=0.5in,right=0.5in ]{geometry}
\usepackage{../styles/ntalgo}
\usepackage[backend=biber]{biblatex}
\addbibresource{refs.bib}

\begin{document}

\title{Entropy Compression}
\author{Navid Tajkhorshid}

\maketitle

Notes based on~\cite{Moser08} and Annie Zheng's Tea Time talk.

\section{The Problem and Algorithm}

Unfortunately, we often don't just want to argue for the \textit{existence} of a
good outcome, but that we can find one ourselves. A canonical example is the
following problem:

\begin{framed}
  \textbf{Low-dependence $k$-SAT}: The input a CNF formula with $n$ variables,
  with each clause containing exactly $k$ variables. It's guaranteed that a
  clause shares variables with at most $d - 1$ other clauses. We want to find a
  satisfying assignment.
\end{framed}

\def\Fix{\textsc{Fix}}
One common way the hypothesis occurs is when each variable appears in at most
$d/k$ clauses. We'll solve this with a recursive randomized ``whack-a-mole''
algorithm. Starting with a random assignment, we'll repeatedly fix unsatisfied
clauses. For each clause $C_i$, denote its dependent clauses as $C_{i, 1}$
through $C_{i, d_i}$ where $d_i \leq d$. Let $C_{i, 0} = C_i$. Also let
$x_{i, 1}, \dots, x_{i, k}$ denote the variables in $C_i$. The following
algorithm fixes an unsatisfied clause while ensuring that new clauses become
unsatisfied.
\begin{centeralgo}
  \Fix($C_i$):\\
  \> pick new unifornmly random values for $x_{i, 1}$ through $x_{i, k}$ \\
  \> while some $C_{i, j}$ is unsatisfied \\
  \> \> $j^* \gets \min\Set{j \mid C_{i, j} \text{ is unsatisfied }}$ \\
  \> \> \Fix($C_{i, j^*}$) \\
\end{centeralgo}
To find a satisfying assignment, we run \Fix over every clause that's
unsatisfied:
\begin{centeralgo}
  \textsc{SatisfyingAssignment}($C_1, \dots, C_m$):\\
  \> pick random values for $x_1, \dots, x_n$ \\
  \> for $i \gets 1$ through $m$ \\
  \> \> if $C_i$ is unsatisfied, \Fix$(C_i)$
\end{centeralgo}
It's not hard to see that if the above algorithm terminates, we have found a
satisfying assignment. The trickiness is in proving termination. In the first
step, each clause is satisfied with probability $(1 - 2^{-k})$, but it's not
clear how it directly helps us.

\section{Entropy Compression}

Let's take a step back. Suppose we want to argue that some randomized algorithm
succeeds with high probability. Let $R$ denote the set of random choices the
algorithm can make; for simplicity assume that each $r \in R$ is equally likely.
We also have a set $F \subseteq R$, which is the set of choices that lead to the
failure. The probability of success is then $1 - \frac{|F|}{|C|}$, and arguing
that the algorithm succeeds with high probability comes down to showing $F$ is
small.

As a simple example, consider the (contrived) problem of finding a small element
in a huge array, where ``small'' means ``in the bottom $1/k$ fraction of
elelemts'' for some $k$. The twist is that the array is far too large to be
scanned through in its entirety. Instead, we sample $s = k\ln n$ elements
uniformly at random, and return the minimum element. The set of random choices
is the set of $s$-length lists of indices, so $|R| = n^s$. We fail iff every
element sampled is not in the bottom $1/k$ fraction, so the set of failing
choices is of size $(n(1 - 1/k))^s$. The probability of failure is
\begin{align*}
  \frac{|F|}{|R|} = \frac{(n(1 - 1/k))^s}{n^s}
                  = (1 - 1/k)^s
                  = ((1 - 1/k)^k)^{\ln n}
                  < 1/n.
\end{align*}

This perspective might seem tedious, particularly when many analyses sidestep
counting sizes of sets and instead work directly with probabilities, but becomes
much more appealing when we take a more information-theoretic approach. Suppose
each failure can be uniquely \emph{encoded} using some small number (say $b$) of
bits. Then $|F|$ cannot be larger than $2^b$. If the set of random choices is
much larger than $2^b$, we can easily argue for success.

Recall the \textsc{SatisfyingAssignment} algorithm described in the last
section. It's easier to think of it as a Monte Carlo algorithm: we set some
threshold $\ell$, and if the total number of calls to $\Fix$ reaches $\ell$
before the formula is satisfied, we report failure. Viewed this way, we can
imagine all the randomness generated upfront: At the beginning, $c = n + k\ell$
random bits are generated (the first $n$ for the initial assignment, the $k\ell$
for the $\ell$ calls to \Fix). Our set $R$ is simply $\Set{0, 1}^{n+k\ell}$--the
set of all bitstrings that could be initially generated.

\subsection{Encoding Failures}

Suppose the algorithm fails on a bitstring $b = b_1\ldots b_c$ (recall that
these are the random choices made upfront). We want to encode this failure in a
small number of bits. Consider the call graph of the algorithm. It's a forest,
where each root is a top-level call to \Fix. Suppose we labeled the top-level
nodes with the index of their clauses, and labeled each non-root node with its
index in its parent's dependency list (that is, the $j$ in $C_{i, j}$). We claim
that with this ordered forest, in addition to the final (failed) assignment, we
can reconstruct the bitstring $b$.

First, replace each non-root node with its global index. Next, work backwards
from the final assignment. If \Fix$(C_i)$ was last called, then the last $k$
bits of $b$ are the assignments to $x_{i, 1}, \dots, x_{i, k}$. Moreover, the
previous assignment must have left $C_i$ unsatisfied, which fully determines the
variables' previous state. We can iterate this process to step back in time,
determining $k$ bits of $b$ at a time, until we reach the initial random
assignment (the first $n$ bits of $b$).

Since this ``call forest'' and final assignment uniquely determine the failed
bistring, we just have to show that they can be encoded efficiently. The final
failed assignment takes $n$ bits to encode. The interesting part is the forest
itself. Let's focus on a single top-level call to \Fix. The tree can be encoded
as an S-expression: any subtree is written as
\begin{align*}
  < \text{subroot} > ( < \text{subtree 1} > \dots).
\end{align*}
For example $1 (2 (3 () 4 () 5 ()) 6)$ uniquely encodes the following labeled
tree: \Tree [.1 [.2 3 4 5 ] 6 ].

In fact, we can save space by omitting each opening parenthesis:
$1 2 3 ) 4 ) 5 )) 6)$. An alternative way to arrive at this encoding is to run a
DFS starting from the root of the tree, writing a ``)'' every time we finish
exploring a subtree. If the tree has $r$ nodes, how many bits do we need? We
need $\ceil{\log_2{m}}$ bits for the label of the root node. We encode ``(''
with a single $1$. Each non-root node label is encoded with $\ceil{\log_2 d
} + 1$ bits by writing a $0$ followed by the index of the neighbor
visited (remember that we don't write the \emph{global} index, instead the index
with respect to the parent recursive call).

\begin{minipage}{0.4\linewidth}
  \Tree [.1 3 1 [ .2 0 ] [ .0 [ .2 3 1 ] ] ].
\end{minipage}
\begin{minipage}{0.5\linewidth}
  \begin{tabbing}
    % CR ntajkhkorshid: Fix the missing parens below here
    1\qquad ( \=3 \quad ( \=) \=1 \quad ( \=) \=2 \quad ( \=0 \quad \=) \=0
    \quad ( \=2 \quad ( \=3 \quad ( \=1 \quad \=) \=) \,) \,) \\
    1 \>3 \>) \>1 \>) \>2 \>0 \>) \>0 \>2 \>3 \>1 \>) \,) \,) \,) \\
    0001 \>011 \>1 \>001 \>1 \>010 \>000 \>1 \>000 \>010 \>011 \>001 \>1 \>1 1 1
    \\
    00010111001101000010000100110011111
  \end{tabbing}
\end{minipage}
\begin{center}
  \small The evolution of the encoding, from the tree to the S-expression to the
  bitstring.
\end{center}

This sums up to a hair under $\ceil{\log_2{m}} + r(2 + \ceil{\log_2 d})$ bits.
To encode the entire forest, we concatenate each tree (note that the tree
encoding allows us to detect when we've finished reading it). Using $\ell$ total
nodes, the forest's encoding uses less than
\[
  m\ceil{\log_2{m}} + \ell(2 + \ceil{\log_2{d + 1}})
\]
bits.

We then have the following:
\begin{lemma}
  A failure after $\ell$ calls to \Fix can be uniquely encoded with at most
  \[
    n + m\ceil{\log_2{m}} + \ell(2 + \ceil{\log_2{d + 1}})
    \leq n + m(\log_2{m} + 1) + \ell(3 + \log_2 d)
  \]
  bits.
\end{lemma}
Call the final quantity $L(\ell)$, the length of the failure encoding if we
allow $\ell$ calls to \Fix. The term to pay attention to is the final one:
$\ell(3 + \log_2 d)$. If $3 + \log_2 d < k \Leftrightarrow d < 2^{k-3}$, the
length of the failure encoding grows slower than the number of random bits
consumed: $n + \ell k$. There can be at most $L(\ell)$ failures, so the
probability of failing is then
\[
  \frac{|F|}{|R|} = \frac{2^{L(\ell)}}{2^{n+\ell k}}
                  = 2^{L(\ell)-(n+\ell k)}
                  = 2^{m(\log_2 m+1)-\ell(k-(3+\log_2 d))}.
\]
For concreteness let's say $d \leq 2^{k-4}$. The probability becomes
$2^{m(\log_2 m+1)-\ell}$, so setting $\ell$ to $O(m \log m + \log n)$ yields a
high probability guarantee. We only need to run \Fix a near-linear number of
times to satisfy the formula!

\paragraph{Remarks.}

It's important to understand what happened. We leveraged the fact that non-top
level calls to \Fix could be encoded with much fewer marginal bits. The ability
to work backwards was essential: no two distinct failing bitstrings could map to
the same encoding. Contrast this with what happens with two successful
bitstrings. They can very well map to the same encoding, since not all the
randomness was actually consumed by the algorithm (if their first $n + 5k$ bits
are common and the algorithm terminates in $5$ steps, then they'll produce the
same solution in the same way).

\end{document}
